\section{Metodologia}
\label{sec:metodologia}

\subsection{Dados}

Gravamos dois vídeos com o celular apoiado, no mesmo quarto:
\texttt{video\_claro} (32\,s, luz acesa) e \texttt{video\_escuro}
(30\,s, pouca luz). Os dois começam com 2 a 3 segundos de cena vazia;
depois a pessoa entra, senta, acena, levanta e chega perto da câmera,
com um ventilador preto, uma cadeira e uma cama ao fundo. A
demonstração ao vivo usa a webcam do notebook ($640 \times 480$).

\subsection{Pipeline}

A Figura~\ref{fig:pipeline} mostra o caminho de cada quadro até a
resposta. Na fase 1, os 60 primeiros quadros, com a cena vazia, treinam
o MOG2 com \texttt{learningRate} $= -1$; na fase 2, cada quadro passa
pelas etapas 0 a 7 com o modelo congelado. Na webcam, uma contagem
regressiva de 5 segundos antes da fase 1 dá tempo de a pessoa sair da
cena.

\begin{figure}[H]
\centering
\begin{tikzpicture}[node distance=5mm and 3mm,
  passo/.style={stepbox, minimum height=1.55cm}]
\node[passo] (e) {\textbf{Entrada}\\quadro da webcam ou do vídeo, altura 480};
\node[passo, right=of e] (p0) {\textbf{0. Compensação}\\ganho por cor = moda de fundo/quadro};
\node[passo, right=of p0] (p1) {\textbf{1. Filtro gaussiano}\\$5 \times 5$};
\node[passo, right=of p1, fill=orange!20] (p2) {\textbf{2. MOG2}\\aprende 60 quadros e depois congela};
\node[passo, right=of p2] (p3) {\textbf{3. Limiar 200}\\descarta a sombra (127)};
\node[passo, below=of p3] (p4) {\textbf{4. Abertura e fechamento}\\elipse $7 \times 7$};
\node[passo, left=of p4] (p5) {\textbf{5. Contornos}\\área $\geq 2000$, preenchidos; contagem};
\node[passo, left=of p5] (p6) {\textbf{6. Buracos}\\maiores que 1500\,px voltam a ser fundo};
\node[passo, left=of p6] (p7) {\textbf{7. Composição}\\erosão $5 \times 5$, desfoque $9 \times 9$, alfa};
\node[passo, left=of p7] (s) {\textbf{Saída}\\original, máscara, resultado e nº de objetos};
\draw[arrow] (e) -- (p0);
\draw[arrow] (p0) -- (p1);
\draw[arrow] (p1) -- (p2);
\draw[arrow] (p2) -- (p3);
\draw[arrow] (p3) -- (p4);
\draw[arrow] (p4) -- (p5);
\draw[arrow] (p5) -- (p6);
\draw[arrow] (p6) -- (p7);
\draw[arrow] (p7) -- (s);
\end{tikzpicture}
\caption{Pipeline do ClearCut-MOG2. Em laranja, o conceito novo.}
\label{fig:pipeline}
\end{figure}

As etapas 1, 3, 4 e 5 são as técnicas da Seção~\ref{sec:tecnicas}. A
etapa 6 existe porque o preenchimento da etapa 5 tapa \textbf{todos} os
buracos, inclusive o vão entre o braço e o corpo com a mão na cintura,
onde reaparecia a parede antiga. Na etapa 7, a máscara é encolhida
antes do desfoque, para não deixar um halo da parede antiga em volta da
pessoa.

\subsection{Parâmetros}

Os parâmetros são constantes comentadas no topo do código
(Tabela~\ref{tab:parametros}). O processamento leva de 27 a 30\,ms por
quadro na webcam, cerca de 35 quadros por segundo, acima dos 30 que ela
entrega.

\begin{table}[H]
\centering
\caption{Parâmetros do pipeline e como foram escolhidos.}
\label{tab:parametros}
\small\setstretch{1.0}
\begin{tabularx}{\linewidth}{>{\raggedright\arraybackslash}p{3.6cm} l >{\raggedright\arraybackslash}X}
\toprule
Constante (etapa) & Valor & Como chegamos \\
\midrule
\texttt{BLUR} (1) & 5 & Valor usado em aula; não variamos. \\
\texttt{FRAMES\_APRENDER} (2) & 60 & Cerca de 2\,s; cabe no início vazio dos vídeos. \\
\texttt{VAR\_THRESHOLD} (2) & 40 & Testamos de 10 a 200. Acima de 40 aumentam os buracos; abaixo, os vídeos melhoram, mas surgem manchas falsas na webcam (Seção~\ref{sec:discussao}). \\
\texttt{LIMIAR} (3) & 200 & Qualquer valor entre 128 e 254 separa sombra (127) de frente (255). \\
\texttt{KERNEL} (4) & 7 & Com 11, o IoU médio caiu de 0,791 para 0,764: a silhueta engorda e perde dedos. \\
\texttt{AREA\_MIN} (5) & 2000 & Maior que as manchas de ruído observadas; não variamos. \\
\texttt{BURACO\_MAX} (6) & 1500 & Na webcam, os vãos entre braço e corpo mediram 2345 e 3104\,px; as falhas de detecção, no máximo 1071\,px. \\
\texttt{ENCOLHER\_BORDA}, \texttt{SUAVIZAR\_BORDA} (7) & 5 e 9 & Comparação visual: sem erosão e com desfoque 15, aparecia um halo claro em volta da pessoa. \\
\texttt{ESPERA\_WEBCAM} & 5\,s & Tempo para sair da cena e para a webcam estabilizar o brilho. \\
\bottomrule
\end{tabularx}
\end{table}

\subsection{Exposição automática: compensação e trava}
\label{sec:exposicao}

O maior problema do projeto veio da câmera, não do MOG2. Quando a
pessoa de camiseta escura entra, a câmera aumenta a exposição sozinha.
Numa gravação de teste, a parede clareou 22 níveis (de 152 para 174),
muito além dos 7,3 que ela tolera (Figura~\ref{fig:mog2}); no vídeo
claro, 92\% da imagem foi marcada como pessoa
(Figura~\ref{fig:exposicao}b).

\textbf{Compensação (etapa 0).} Multiplicamos cada canal do quadro por
um ganho que o devolve ao brilho do fundo aprendido. O ganho é o valor
\textbf{mais comum} (a moda do histograma) da razão fundo
aprendido / quadro: quando a câmera clareia, todos os pixels do fundo
pedem o mesmo ganho e formam um pico, enquanto os da pessoa ficam
espalhados; por isso funciona mesmo com a pessoa ocupando mais da
metade da imagem (Figura~\ref{fig:exposicao}c). Foi a terceira versão:
a mediana da imagem inteira era puxada pela pessoa, e a mediana só do
fundo do quadro anterior chegava a inverter a máscara.

\textbf{Trava (webcam).} Na webcam, com a pessoa perto, a parede foi de
130 para 235, e de 19\% a 31\% dos pixels estouraram em 255. A câmera
também muda a curva de tons: no mesmo quadro, a parte clara da parede
pedia ganho 0,70 e a escura, 0,59, e nenhum ganho único corrige as duas
(Figura~\ref{fig:exposicao}d--e). Por isso, depois da contagem, o
programa \textbf{trava} a exposição pelo OpenCV
(\texttt{CAP\_PROP\_AUTO\_EXPOSURE}), testando cinco valores e ficando
com o de brilho mais próximo do automático. Ao reaprender o fundo
(tecla B) ou sair, a câmera volta ao automático. Com a trava, a parede
ficou entre 108 e 118 o teste inteiro (Figura~\ref{fig:exposicao}f--g).

\begin{figure}[H]
\centering
\includegraphics[width=\linewidth]{fig_exposicao.jpg}
\caption{Exposição automática. Vídeo claro: (a) quadro, (b) máscara sem
compensação e (c) com compensação. Webcam: (d--e) sem a trava, o fundo
some; (f--g) com a trava, a máscara mostra só a pessoa.}
\label{fig:exposicao}
\end{figure}

\subsection{Avaliação}

Escolhemos 12 quadros, 6 de cada vídeo, todos com a pessoa e misturando
casos fáceis e difíceis. Quadros vazios ficaram de fora, porque programa
e gabarito vazios dão IoU $= 1$ e inflariam o resultado. Para cada
quadro, o grupo pintou à mão a máscara correta (gabarito) com uma
ferramenta própria (\texttt{criar\_gabarito.py}), e o
\texttt{avaliacao.py} calcula o IoU e conta \textbf{acerto} quando
IoU $\geq 0{,}8$. Também conferimos a contagem de objetos nos 78 quadros
salvos (1 a cada 20). Como o sistema não classifica objetos em classes,
não há matriz de confusão.
